% ===========================================================
%  THEOREM 1 -- FORMAL SETUP
%  Direction-Invariant Label Flip under Linear Batch Correction
%  (v10 AAAI 2027)
% ===========================================================

\documentclass[11pt,a4paper]{article}
\usepackage[margin=1in]{geometry}
\usepackage{amsmath,amssymb,amsthm,bm,hyperref,booktabs}
\newtheorem{theorem}{Theorem}
\newtheorem{corollary}[theorem]{Corollary}
\newtheorem{proposition}[theorem]{Proposition}
\newtheorem{lemma}[theorem]{Lemma}
\newtheorem{assumption}{Assumption}
\newtheorem{definition}{Definition}
\title{Theorem 1 -- Formal Setup\\
 \large Direction-Invariant Label Flip under Linear Batch Correction}
\date{v10 AAAI 2027 -- 2026-04-24}
\begin{document}
\maketitle

\section{Objects and notation}

Let
\begin{align*}
  X &\in \mathbb{R}^{n\times p}          &&\text{row-wise expression matrix,}\\
  Y &\in \{0,1\}^{n}                     &&\text{binary clinical label,}\\
  B &\in \{1,\dots,K\}^{n}               &&\text{cohort (batch) assignment, }K\ge 2,\\
  \pi_{k} &= \{\,i : B_i=k\,\}           &&\text{index set of samples in batch }k.
\end{align*}
All rows $X_i$ are assumed column-standardised across the full dataset
so that sample-level covariance is order-one.

\subsection{Key directions}

We define the three unit vectors that drive the analysis.
\begin{align}
  w_Y \;&=\; \frac{\,\mu_1-\mu_0\,}{\lVert\mu_1-\mu_0\rVert},
     &\mu_c &=\mathbb{E}[X_i\mid Y_i=c] \quad (c\in\{0,1\}),
     \tag{label direction}\\[2pt]
  w_B \;&=\;
     \underset{\lVert w\rVert=1,\; w\perp \ker(P_Y^{\perp})}{\arg\max}\;
     w^{\!\top}\,\mathrm{Cov}_{\!B}(X)\,w,
     &\mathrm{Cov}_{\!B}(X) &=
     \mathbb{E}_{k}\!\bigl[(\mu_{B=k}-\bar\mu)(\mu_{B=k}-\bar\mu)^{\!\top}\bigr],
     \tag{dominant batch direction}\\[2pt]
  w_{\text{trained}} \;&=\;
     \arg\min_{w,b} \sum_{i=1}^{n}\ell\!\bigl(Y_i,w^{\!\top}T_B(X)_i+b\bigr)
     + \lambda\lVert w\rVert^{2},
     \tag{trained classifier}
\end{align}
where $\ell$ is a convex margin loss (logistic, hinge, or squared-hinge)
and $\lambda\ge 0$ is a ridge penalty.

We write $\rho = \langle w_Y, w_B\rangle \in [-1,1]$ and note that under
the sign convention above $\rho\ge 0$ after at most a relabelling of the
batch coordinate.

\subsection{Projection operators}

\begin{definition}[Batch subspace projector]
  $P_B \;=\; U_B U_B^{\!\top}$, where $U_B\in\mathbb{R}^{p\times r}$
  has orthonormal columns spanning the column space of the centred
  batch-indicator matrix $\tilde B\in\mathbb{R}^{n\times K}$ acting on
  $X$; i.e.\ $\mathrm{range}(P_B)=\mathrm{span}\{\mu_{B=k}-\bar\mu:k=1,\dots,K\}$.
  Rank $r = K-1$ generically.
\end{definition}

\begin{definition}[Label direction projector]
  $P_Y \;=\; w_Y w_Y^{\!\top}$.
\end{definition}

\subsection{Linear batch correction with label covariate}

Following the standard ComBat-with-covariate formulation, we model
\begin{equation}
   X_i \;=\; \alpha + \beta^{\!\top} Y_i + \gamma_{B_i} + \varepsilon_i,
   \qquad \varepsilon_i\sim\mathcal{N}(0,\Sigma).
   \label{eq:additive-model}
\end{equation}
The linear correction operator is
\begin{equation}
  T_B(X_i) \;=\; X_i - \hat\gamma_{B_i}
  \;=\; X_i - \bigl(\bar X_{\pi_{B_i}} - \bar X - \hat\beta^{\!\top}(\bar Y_{\pi_{B_i}}-\bar Y)\bigr),
  \label{eq:TB}
\end{equation}
where $\hat\gamma_k$ is the OLS estimate of $\gamma_k$ holding the label
covariate $Y$ fixed. Matrix form: there exists $M_B\in\mathbb{R}^{n\times n}$
(idempotent on the batch-indicator column space, orthogonal to~$Y$) such that
\begin{equation}
  T_B(X) \;=\; (I-M_B)X \;+\; (\text{label-covariate adjustment}).
  \label{eq:TB-matrix}
\end{equation}

\subsection{The key quantity}

The \emph{alignment parameter}
\begin{equation}
  \rho^{2} \;=\; \bigl\langle w_Y,\,w_B\bigr\rangle^{2}
  \;\in\;[0,1]
  \label{eq:rho}
\end{equation}
measures how closely the dominant batch direction parallels the true
label direction.  \textbf{Theorem 1} (stated below) shows that
$\rho^2>\tfrac12$ is the sharp threshold at which the trained classifier
under $T_B$ \emph{inverts} relative to the pre-correction discriminant.

\section{Assumptions}

\begin{assumption}[Additive generative model]
\label{asm:additive}
Equation~\eqref{eq:additive-model} holds; residuals are isotropic with
sub-Gaussian tails.
\end{assumption}

\begin{assumption}[Well-posedness]
\label{asm:wellposed}
$\mu_1\ne\mu_0$ and at least two batches differ in mean; the ridge
parameter satisfies $\lambda = o(n)$.
\end{assumption}

\begin{assumption}[Cohort imbalance or near-confounding]
\label{asm:conf}
There exist $k\ne k'$ with $\bigl|\Pr(Y=1\mid B=k)-\Pr(Y=1\mid B=k')\bigr|\;\ge\;\delta>0$.
This is what makes $\rho$ strictly positive in the finite-sample regime.
\end{assumption}

\begin{assumption}[Sample scaling]
\label{asm:n}
$n/p$ is bounded away from zero as $n\to\infty$; $\min_k|\pi_k| = \Omega(n/K)$.
\end{assumption}

\section{Main theorem (statement)}

\begin{theorem}[Direction-invariant label flip]
\label{thm:flip}
Assume~\ref{asm:additive}--\ref{asm:n}. Let $f_{\text{trained}}(x)=
w_{\text{trained}}^{\!\top}x + b_{\text{trained}}$ be the linear
classifier fit on $(T_B(X),Y)$. Let
\[
   \mathrm{AUC}_{\text{post}}
   \;=\; \Pr\bigl(f_{\text{trained}}(X^\ast)>f_{\text{trained}}(X^\dagger)\,\bigm|\, Y^\ast=1,\,Y^\dagger=0\bigr)
\]
on a held-out cohort, and set
\[
   \mathrm{DIAL}
   \;=\; \max\!\bigl(\mathrm{AUC}_{\text{post}},\;1-\mathrm{AUC}_{\text{post}}\bigr)
         \;-\; \mathrm{AUC}_{\text{post}}.
\]
Then there exists a universal threshold $\rho^\star\in(0,1)$ with
$\rho^{\star 2}=\tfrac12$ such that
\begin{enumerate}
\item[\textup{(i)}] if $\rho^2<\rho^{\star 2}$,
 $\lim_{n\to\infty}\Pr(\mathrm{AUC}_{\text{post}}>1/2)=1$,
 and $\mathrm{DIAL}\xrightarrow{p}0$;
\item[\textup{(ii)}] if $\rho^2>\rho^{\star 2}$ and Assumption~\ref{asm:conf} holds,
 $\lim_{n\to\infty}\Pr(\mathrm{AUC}_{\text{post}}<1/2)=1$,
 and $\mathrm{DIAL}\xrightarrow{p} 1-2\,\mathrm{AUC}_{\text{post}}^{\infty}>0$.
\end{enumerate}
Equivalently, $\mathrm{sign}\langle w_{\text{trained}},w_Y\rangle$ flips
sharply at $\rho^2=\tfrac12$.
\end{theorem}

\begin{corollary}[DIAL characterisation]
\label{cor:dial}
Under Theorem~\ref{thm:flip}, $\mathrm{DIAL}>0$ if and only if
$\rho^2>\tfrac12$.  In particular, DIAL is a direct estimator of the
mass on the $\rho^2>\tfrac12$ event conditional on cohort structure.
\end{corollary}

% Proof is in theorem1_proof.tex.

\end{document}
